\documentclass{article}
\usepackage[margin=0.5in]{geometry}
\usepackage{graphicx}
\usepackage{caption}
\usepackage{amsmath}

\title{Winterizing the Rainwater System}
\date{2025 May}
\author{JC Vaught, William Wenzel, Tanner Harmon, Presston McConnell}

\begin{document}
\maketitle

\section*{Overview}
This guide is intended for use when sub-\(20^\circ\mathrm{F}\) (approx.\ \(-6^\circ\mathrm{C}\)) temperatures are expected for multiple consecutive days and the system will not be in service. A preferable winterization technique is to place a 60\,W incandescent bulb in the pump house, powered via an extension cord, to help maintain temperatures above freezing instead of draining the system.

\section*{Procedure}

\subsection*{1. Turn Off the Inverter}
Turn off the inverter. The green LED indicator will extinguish, confirming that the unit is powered down.

\begin{figure}[h]
  \centering
  \includegraphics[width=0.55\linewidth]{Work_Products/inverter.png}
  \caption{Verify that the inverter is turned off.}
  \label{fig:inverter-off}
\end{figure}

\newpage\subsection*{2. Close the Main Valve}
Rotate the main valve handle to the closed position (lever perpendicular to the flow path).

\begin{figure}[h]
  \centering
  \includegraphics[width=0.7\linewidth]{Work_Products/main_valve.png}
  \caption{Main valve in the closed position.}
  \label{fig:main-valve}
\end{figure}

\subsection*{3. Drain the System}
Attach a garden hose to the spigot inside the pump house. Open the spigot and allow all water to drain until flow ceases.

\begin{figure}[h]
  \centering
  \includegraphics[width=0.7\linewidth]{Work_Products/spigot.png}
  \caption{Draining the system via the pump-house spigot.}
  \label{fig:draining}
\end{figure}

Once the hose is disconnected and no more water flows, close the spigot to complete the procedure.

\end{document}
